金融社交媒体上的投资者评论蕴含丰富的情绪信息，准确识别其情感倾向是金融舆情监测和投资者情绪研究的基础。大语言模型在文本理解上表现突出，LoRA 等参数高效微调方法使其能够在单张消费级显卡上完成微调，但其在中文金融情感分类上相对传统方法的优势是否可靠、给出的置信度能否被信任，需要在明确数据质量与评估口径后加以判断。

本文以中文金融社交媒体情感三分类数据集 FinFE 为对象，在单张 RTX 3090 显卡上比较了 TF-IDF 配合逻辑回归和支持向量机、BERT 与 RoBERTa-wwm-ext 微调，以及 Qwen3-8B、Qwen2.5-7B-Instruct、GLM-4-9B 三个开源大模型的零样本、少样本提示与 LoRA 微调，并从统计推断、置信度可靠性和数据质量三个角度进行分析。

研究发现：（1）原始评估文件中约 54\% 的记录与原始训练文件存在规范化文本重叠；本文区分训练与验证曝光，清洗后得到集合间没有完全重复文本的数据划分，并在共同未曝光样本上比较清洗策略。（2）LoRA 微调的大模型宏平均 F1 为 0.798～0.807，高于 RoBERTa 的 0.778 和 TF-IDF 的 0.70；零样本和少样本提示不及有监督微调。（3）在固定数据划分的三个训练种子中，大模型 F1 点估计均更高，多种子方差分析的事后比较显著；但种子 42 的 McNemar 检验和配对 Bootstrap 经多重比较调整后，未检出大模型相对 RoBERTa 的显著优势。（4）Qwen3 零样本平均置信度为 96.4\%，准确率为 67.1\%；LoRA 模型经测试分布内五折温度交叉拟合后，期望校准误差为 0.013～0.030，只回答高置信度的一半样本时准确率约为 96\%。标准保形预测的整体覆盖率约为 90\%，约 77\%～78\% 的集合为单元素，但该子集准确率为 88.7\%～89.3\%，不具有单独的 90\% 正确率保证。（5）重复记录的配对一致率为 82.9\%，82\% 的冲突涉及中性类；由于缺少独立复标证据，不将噪声模型推算值解释为实测准确率上限。种子 42 的错误分析中，Qwen3 LoRA 的 307 条错误有 250 条涉及中性类；固定抽样案例用于说明边界，不作因果归因。此外，单次 4 比特 QLoRA 实验在框架记录的峰值分配显存 11.2 GB 下取得了与 16 位 LoRA 接近的效果。

本文的结果表明，在评估大模型时应同时关注数据质量、统计显著性和置信度的可靠性，可为金融机构开展舆情监测时的模型选型与人机协作提供参考。

\textbf{关键词：} 大语言模型；LoRA 微调；金融情感分析；统计检验；置信度校准
